\scanpage{42}
\mspage{16}

ne donne quelque chose que si on la complète par la prescription que l'on prend une métrique ``distinguée'' sur $X^{\mathrm{or}}$ qui soit stable par $\sigma$.

Considérons maintenant les seuls cas nouveaux -- ceux où $X$ est 0-connexe et $\partial X\ne\varnothing$, correspondant à une $X^{\mathrm{or}}=\mathscr X$ connexe, munie d'une antiinvolution $\sigma$ telle que $\deleted{\unsure{\mathbb C}}\mathscr X^\sigma\ne\varnothing$. Les seuls cas problématiques pour l'existence d'une métrique canonique sur $X$, sont les suivants :

a') $\mathscr X\simeq\mathbb C$. \deleted{\unsure{Les anti.|Les aut.}}\inserted{OPS} $\mathscr X=\mathbb C$. Les anti-aut. de $\mathbb C$ sont
\[
z\longmapsto a\bar z+b\qquad(a\in\mathbb C^{\times},b\in\mathbb C)
\]
Si on prend une antiinvolution à pts fixes, OPS que $0$ est pt fixe, donc $b=0$, donc \deleted{l'application est $z\longmapsto a\bar z$}
\[
\sigma z=\deleted{\unsure{\alpha|a}}a\bar z\qquad\text{avec}\quad a\bar a=1\quad\text{i.e}\quad |a|=1\quad\text{i.e}\quad a\in\mathbb U
\]
(pour avoir $\sigma^2=\mathrm{id}$) soit

\deleted{l'ens des pts fixes est}
\[
\sigma z=\alpha^2\bar z=\alpha/\bar\alpha\,\bar z\qquad\text{si}\quad a=\alpha^2,\quad a\in\mathbb U,
\]
et l'ens des pts fixes est
\[
\mathbb C^\sigma=\alpha.\mathbb R
\]
\margin{symétrie p.r. à la droite $\alpha.\mathbb R$}
% archive sheet number: 44
